\documentclass{article}

% ============================================================
% MA 213 In-Class Worksheet Template
% Student-facing worksheet for lecture activities.
% ============================================================

\usepackage[margin=1in]{geometry}
\usepackage{amsmath}
\usepackage{xcolor}
\usepackage{parskip}
\usepackage{array}
\usepackage{enumitem}
\usepackage{booktabs}

\newcommand{\coursenumber}{MA 213}
\newcommand{\weekplaceholder}{10}
\newcommand{\lectureplaceholder}{24}
\newcommand{\lecturetitle}{Applying the CLT for the Sample Mean: Think/Pair/Share}

\newcommand{\nameblank}{\rule{1.75in}{0.4pt}}
\newcommand{\idblank}{\rule{1.55in}{0.4pt}}
\newcommand{\shortblank}{\rule{1in}{0.4pt}}
\newcommand{\longblank}{\rule{0.93\linewidth}{0.4pt}}

\begin{document}

\begin{center}
  {\LARGE \coursenumber{} Week \weekplaceholder{}: Lecture \lectureplaceholder{}}\\
  {\large \lecturetitle{}}
\end{center}

\begin{enumerate}[label=\arabic*.,leftmargin=*,itemsep=.7em]
  \item \textbf{Your name:} \nameblank \hfill \textbf{BU ID:} \idblank
\end{enumerate}

\textbf{Big idea:} The CLT for the sample mean tells us $\bar{X} \sim N\left(\text{mean}=\mu, SE = \frac{\sigma}{\sqrt{n}}\right)$, and this holds \emph{exactly} for any $n$ when the population itself is Normal. Let's use it to build a confidence interval by hand.

\section*{The Setup}

You take a random sample of size $n = 10$ from a population and calculate the sample mean $\bar{X} = 49$. Assume that the population is Normal with $\sigma = 20$.

\section*{Step 1: Think (2 mins)}

\textbf{(a)} Using $\bar{X} \sim N\left(\text{mean}=\mu, SE = \frac{\sigma}{\sqrt{n}}\right)$, calculate a 95\% confidence interval for the population mean $\mu$.
\vspace{1.1in}

\textbf{(b)} Is the data consistent with a population mean of 60? Why or why not?
\vspace{0.6in}

\newpage
\section*{Step 2: Pair (2 mins)}

\begin{enumerate}[label=\arabic*.,leftmargin=*,itemsep=.7em, start=2]
  \item \textbf{Partner's name:} \nameblank \hfill \textbf{BU ID:} \idblank
\end{enumerate}

Compare your confidence interval and conclusion with your partner. Then discuss:

\begin{itemize}
  \item Did you both use the same critical value $z^\star$? Where does it come from?
  \item Would your approach change if the population were \emph{not} known to be Normal, but $n$ were still only 10? What if $n$ were 100 instead?
  \item Would your approach change if $\sigma$ were unknown instead of $\sigma = 20$?
\end{itemize}

\textbf{Our thoughts:}
\vfill

\end{document}
